\scanpage{10}
\mspage{9}

\begin{center}
CHAPITRE VIII

CONS\unsure{T}RUCTION de MESURES INVARIANTES

§ 1. \uline{Construction d'une mesure de Haar.}
\end{center}

\uline{Dans ce} §, $G$ \uline{désignera un groupe localement compact}.

1. \uline{Le théorème d'existence et d'unicité}.

Rappelons (Chap.VII, § 2, n°4) qu'on appelle \uline{mesure de Haar} sur $G$ toute mesure positive $\alpha\neq0$ sur $G$, invariante à gauche, et qu'une telle mesure est alors une jauge propre à gauche et à droite, satisfaisant à une formule de la forme
\[
d\alpha(axb)=\Delta_G(b)d\alpha(x)
\]
quels que soient $a\in G$, $b\in G$, $\Delta_G$ étant un caractère de $G$ ; on a alors
\[
d\alpha(x^{-1})=\Delta_G(x^{-1})d\alpha(x).
\]
Si $\Delta_G=1$, $\alpha$ est invariante à gauche et à droite ; on dit que $\alpha$ est \uline{bi-invariante} ; il en est évidemment ainsi si $G$ est commutatif.

THEOREME 1.-- \uline{Sur tout groupe localement compact} $G$, \uline{il existe une mesure de Haar, et, à un facteur constant près, il n'en existe qu'une}.

L'unicité a déjà été démontrée (Chap.VII, § 2, n°4, th.1) ; on va démontrer l'existence par une construction explicite. Introduisons d'abord une notation. Soit $\mu=\sum_{i=1}^{n}a_i\varepsilon_{s_i}$ une mesure de support fini $\{s_1,\ldots,s_n\}$ sur $G$ ; alors, pour $f\in\mathcal{K}(G)$, on désignera par $\gamma(\mu)f$ la fonction $f'\in\mathcal{K}(G)$ définie par
\[
f'(x)=\sum_{i=1}^{n}a_i f(s_i^{-1}x)\ ;
\]
\margin{\figure{\begin{tikzpicture}[x=1cm,y=1cm]
\draw (0,2.5) .. controls (0.7,2.65) and (1.3,2.05) .. (0.85,2.2);
\draw (0.15,1) .. controls (1.5,0.8) and (1.35,0.55) .. (0.2,0);
\end{tikzpicture}}}

on notera que, dès qu'on aura construit une mesure de Haar, on aura le droit, en vertu des définitions du Chap.VII, § 2, d'écrire $\mu*f$ au lieu de $\gamma(\mu)f$ (cf. la Remarque à la fin du Chap.VII, § 2, n°2). L'application $(\mu,f)\longrightarrow\gamma(\mu)f$ ainsi définie est évidemment bilinéaire ; comme on a $\gamma(\varepsilon_{st})=\gamma(\varepsilon_s)\gamma(\varepsilon_t)$, on en conclut qu'on a $\gamma(\mu*\nu)=\gamma(\mu)\gamma(\nu)$. De plus, $\gamma(\mu)f$ est $\geq0$ si $\mu$ et $f$ sont $\geq0$ ; et, si $\mu$ et $f$ sont $\geq0$, on ne peut avoir $\gamma(\mu)f=0$ que si $\mu=0$ et $f=0$.
